%  Appendices A-C of the Electronics manuscript (labels app:prompt, app:series, app:tables).
%
%  A is transcribed from the code, not typed.  The prompt is MechanicsPromptBuilder.build_regime_prompt with the 30 data lines
%  replaced by the placeholder {gex_data_table}; the schema and the request are common.REGIME_JSON_SCHEMA and common.chat_body.
%  The three typographic symbols of the prompt (U+2264, U+2265, U+2192) are written in ASCII (<=, >=, ->) in this file and
%  typeset as symbols by the "prompt" listing style, which keeps the sources ASCII.
%  B.1 to B.4 are the four statements of the data audit of 2026-10-28 (sql_series_a_frozen, _b_new, _c_ex0dte and
%  _d_intraday.sql) VERBATIM, comments included; B.5 adds the statement of the scaling factor (tested read-only on 2026-10-29)
%  and describes the recovered series.  Do not edit the SQL lines by hand: change the statement and say so in the text.
%  Re-check the listings against the code whenever the prompt, the schema or the model configuration changes
%  (README.md, section "Keeping Appendices A and B in step with the code").
%  Appendix C is filled from results/ (05_report.py, 2026-10-31); every table carries its source in a "% src:" comment.
%  Prompt check (2026-10-31): the listing of A.2, with the data lines replaced by the placeholder and the three symbols in ASCII,
%  equals common.build_prompt on 34,244 of 34,244 obfuscated gpt-6-luna requests stored in data/backups/gex_llm/*/input_*.jsonl.
\appendixtitles{yes}
\appendix

% Listing styles (the listings package is loaded in the main file).  "prompt" typesets <=, >= and -> as symbols; it is
% used for the prompt only, because the SQL contains the operators literally.  The two penalties are set to 0 inside the
% listing because mdpi.cls loads hyphenat[none], which makes both infinite and so forbids the discretionary breaks that
% listings uses for breaklines (without this the long prompt and SQL lines run into the margin).
\lstdefinestyle{appcode}{basicstyle={\ttfamily\scriptsize\exhyphenpenalty=0 \hyphenpenalty=0 },
  columns=fixed, keepspaces=true, showstringspaces=false, upquote=true,
  breaklines=true, breakindent=0pt, postbreak=\mbox{\textcolor{gray}{$\hookrightarrow$}\space},
  aboveskip=6pt, belowskip=6pt, xleftmargin=4pt, backgroundcolor=\color{gray!8}}
\lstdefinestyle{prompt}{style=appcode, breakatwhitespace=true,
  literate={<=}{{$\leq$}}1 {>=}{{$\geq$}}1 {->}{{$\rightarrow$}}1}

%=================================================================
\section{Prompt, Schema and API Configuration}
\label{app:prompt}

This appendix reproduces the complete request sent to the model for each 30-day window, so that the experiment can be repeated from the manuscript alone. The prompt text is the output of \ident{MechanicsPromptBuilder.build_regime_prompt} (file \path{src/llm/mechanics_prompt_builder.py}); the schema and the request body are the constant \ident{REGIME_JSON_SCHEMA} and the function \ident{chat_body} of \path{scripts/validation/electronics/common.py}. The listings were generated from these objects and compared with them.

\subsection{Model and API Configuration}
\label{app:prompt:config}

% src: scripts/validation/electronics/common.py BATCH_MAX_REQUESTS=1500, MAX_COMPLETION_TOKENS=4000, RETRY_MAX_COMPLETION_TOKENS=8000; cost_ledger.yaml pricing_usd_per_mtok gpt-6-luna batch input 0.05 output 0.25 (standard 0.1/0.5); 02_decisions_log.md line 60 (model created 2026-09-14T17:15:02Z); reasoning effort medium
\begin{itemize}[leftmargin=*,labelsep=5.5mm]
\item \textbf{Model:} \ident{gpt-6-luna}, an OpenAI model alias without a dated snapshot. The creation timestamp of the model object (2026-09-14) is stored with every run and the response identifier with every answer. The published run used \ident{o4-mini-2025-04-16}.
\item \textbf{Access:} OpenAI Batch API, endpoint \path{/v1/chat/completions}, completion window of 24 hours, at most 1{,}500 requests per batch job. Each input line is one request that carries the key \texttt{<arm>|<symbol>|<window\_end>} in \ident{custom_id}; a test-retest duplicate adds the suffix \texttt{\#rep1}, and its arm is stored as \texttt{<arm>\#rep1}.
\item \textbf{Reasoning effort:} \ident{medium}, the default effort of o4-mini in the published run.
\item \textbf{Maximum completion tokens:} 4{,}000; reasoning tokens count against the limit. An answer that ends with \ident{finish_reason} = \ident{length} is requested once more with 8{,}000.
\item \textbf{Response format:} the strict JSON schema of Appendix~\ref{app:prompt:schema}. The earlier runs requested the same nine fields in the prompt text only; this is the one protocol change of the study.
\item \textbf{Not sent:} temperature, seed, system message, tools and stop sequences.
\item \textbf{Parsing:} an answer is parsed as JSON; a second attempt tolerates raw control characters inside strings, and a last fallback extracts the JSON object from free text. The parse status of every answer (\ident{ok}, \ident{fallback}, \ident{failed} or \ident{missing}) is stored.
\item \textbf{Prices:} USD 0.05 per million input tokens and USD 0.25 per million output tokens on Batch, twice that on the standard API.
\end{itemize}

Identical requests need not return identical answers, because no seed is set. The test-retest duplicates of Section~\ref{sec:methods:arms} measure this variation, and the deterministic classifier, not the model, supplies the reference label of each window.

\subsection{Prompt Text}
\label{app:prompt:body}

The prompt is a single user message. The placeholder \texttt{\{gex\_data\_table\}} stands for the 30 data lines of a window, one per day, in the form \texttt{Day T-29: +3.42B} to \texttt{Day T+0: -12.18B}: a relative day label, the net GEX in billions of US dollars with two decimals and an explicit sign, and neither a date nor a ticker. In the two dated arms (\ident{p34_dated_2024} and \ident{p34_dated_2020}) the label is instead the ISO date of the trading day, as in the November 2025 runs (\texttt{2024-01-02: -33.01B}). The rest of the text is the same for every window, series and symbol (prompt version \ident{v2.0_postgresql}) and is the text of the Phase 4A run.

The listing was checked against the stored requests. With the 30 data lines replaced by the placeholder, it equals the user message of 34{,}244 of 34{,}244 requests of the obfuscated gpt-6-luna input files (the dated-label arms differ only in the day labels). Each stored answer also carries the sha256 of the prompt that produced it in \ident{detection_metadata.prompt_sha256}.
% src: comparison of this listing with the user message of every input_*.jsonl of the non-dated gpt-6-luna runs in data/backups/gex_llm/ (34,244 requests, 34,244 equal; the 40-request synchronous pilot is included; 34,712 ledger requests minus 4 x 234 dated-arm requests), 2026-10-31; also equal to common.build_prompt

\begin{lstlisting}[style=prompt]
You are a market structure analyst specializing in dealer gamma positioning regimes.

TASK: Analyze this 30-day period and determine if it represents a PERSISTENT regime where dealer constraints create forced, directional flows.

## 30-DAY GEX DATA

{gex_data_table}

## REGIME CLASSIFICATION FRAMEWORK

### PERSISTENT REGIMES (Detect These)

**1. PERSISTENT POSITIVE REGIME**
- Definition: Dealers are LONG gamma, forced to sell into strength
- Criteria:
  * >70% of days (21+/30) have positive net GEX
  * Average magnitude >$5B
  * <=5 sign flips across 30 days
  * Stable directional constraint

**Mechanism**: When dealers hold long gamma:
- Price rises -> Dealers MUST sell shares (rebalance)
- Price falls -> Dealers MUST buy shares (rebalance)
- Creates dampening, mean-reverting flows
- Constraint is STRUCTURAL (dealers cannot avoid)

**2. PERSISTENT NEGATIVE REGIME**
- Definition: Dealers are SHORT gamma, forced to buy into strength
- Criteria:
  * >70% of days (21+/30) have negative net GEX
  * Average magnitude >$5B
  * <=5 sign flips across 30 days
  * Stable directional constraint

**Mechanism**: When dealers hold short gamma:
- Price rises -> Dealers MUST buy shares (chase)
- Price falls -> Dealers MUST sell shares (chase)
- Creates amplifying, momentum flows
- Constraint is STRUCTURAL (dealers cannot avoid)

---

### NON-REGIMES (Reject These)

**3. TRANSITIONAL (Reject)**
- Frequent sign flips between positive/negative GEX
- No dominant regime direction (less than 70% same sign)
- Market in regime change period
- Example: 15 positive days, 15 negative days (50/50 split)

**Why Reject**: No persistent constraint. Dealers face mixed conditions daily. Not a structural regime.

**4. LOW CONVICTION (Reject)**
- Consistent sign BUT weak magnitude (<$5B average)
- Example: 25 days positive, avg $2B GEX
- Insufficient constraint to create persistent forced flows

**Why Reject**: Even if sign is consistent, magnitude too weak to force dealers into meaningful positions. Not a structural constraint.

---

## ANALYSIS QUESTIONS

Systematically evaluate the 30-day window:

**Step 1: Sign Persistence**
1. Count days with positive net GEX
2. Count days with negative net GEX
3. Calculate persistence percentage: max(positive_days, negative_days) / 30 * 100
4. Does it meet 70% threshold (21+ days)?

**Step 2: Magnitude Assessment**
1. Calculate average GEX magnitude (absolute value): sum(|net_gex|) / 30
2. Is average magnitude >=$5B?
3. Check for extreme outliers that might distort average

**Step 3: Stability Check**
1. Count sign flips: How many times does GEX switch from pos->neg or neg->pos?
2. Are there <=5 sign flips across 30 days?
3. Stable regime should have low flip count

**Step 4: Regime Classification**
- If Steps 1, 2, 3 all pass AND positive dominates -> PERSISTENT POSITIVE
- If Steps 1, 2, 3 all pass AND negative dominates -> PERSISTENT NEGATIVE
- If Step 1 passes but Step 2 fails -> LOW CONVICTION (reject)
- If Step 1 fails -> TRANSITIONAL (reject)

---

## CONFIDENCE CALIBRATION (Mechanical Guidance)

Use these concrete anchors to calibrate confidence:

**90-100 (Very High Confidence)**
- 25-30 days same sign (83-100% persistence)
- Average magnitude >$10B
- 0-2 sign flips (highly stable)
- Example: "29 negative days, avg $15B, 1 flip"

**70-89 (High Confidence)**
- 21-24 days same sign (70-80% persistence)
- Average magnitude $5-10B
- 2-4 sign flips (moderately stable)
- Example: "23 negative days, avg $7B, 3 flips"

**50-69 (Borderline - Use with Caution)**
- 18-20 days same sign (60-67% persistence)
- Average magnitude $3-5B
- 5-7 sign flips
- Example: "20 negative days, avg $4B, 6 flips"
- **Note**: Borderline cases should generally be REJECTED unless other factors strengthen confidence

**0-49 (Reject - Not Persistent)**
- <18 days same sign (<60% persistence)
- OR average magnitude <$3B
- OR >7 sign flips
- These are NOT persistent regimes

**Important**: Confidence is a FILTER, not a probability. Use it to distinguish clear regimes (70+) from borderline (50-69) from noise (<50).

---

## OUTPUT FORMAT (JSON)

Provide your analysis in this exact JSON structure:

```json
{
    "regime_detected": true/false,
    "regime_type": "persistent_positive|persistent_negative|transitional|low_conviction",
    "positive_days": <count as integer>,
    "negative_days": <count as integer>,
    "avg_magnitude_billions": <value as number>,
    "sign_flips": <count as integer>,
    "persistence_pct": <percentage as number>,
    "confidence": <integer 0-100>,
    "reasoning": "Explain step-by-step why this is/isn't a persistent regime. Reference specific metrics (persistence %, avg magnitude, sign flips). If rejecting, state which criterion failed."
}
```

**IMPORTANT**: All numeric fields (confidence, positive_days, negative_days, sign_flips, avg_magnitude_billions, persistence_pct) MUST be numbers (integers or decimals), NOT words like "thirty-five" or "fifty".

**regime_detected Rules**:
- `true` ONLY if regime_type is "persistent_positive" or "persistent_negative"
- `false` if regime_type is "transitional" or "low_conviction"

---

## KEY PRINCIPLES

1. **Selectivity is Expected**: Most windows will NOT be persistent regimes (expect 30-50% detection rate)

2. **ALL Criteria Must Pass**: Persistence + Magnitude + Stability required for detection

3. **Rejection is Valid**: Saying "no persistent regime" is a correct answer for transitional/weak periods

4. **Mechanical Over Qualitative**: Use concrete thresholds (70%, $5B, 5 flips) rather than subjective judgment

5. **Structural Focus**: Only detect when dealers are FORCED into directional positions by constraints

Analyze the 30-day GEX data above and provide your regime classification in JSON format.
\end{lstlisting}

\subsection{Output Schema}
\label{app:prompt:schema}

The answer is constrained by the following strict JSON schema, passed as \ident{response_format}. With \ident{strict} set, the model can neither omit nor add a field. The boolean \ident{regime_detected} is the primary label of the study and \ident{regime_type} the secondary one; the counts, the mean magnitude and the persistence are the model's own report and are compared with the values of the deterministic classifier (Section~\ref{sec:results:fidelity}).

\begin{lstlisting}[style=appcode]
{
  "type": "json_schema",
  "json_schema": {
    "name": "regime_classification",
    "strict": true,
    "schema": {
      "type": "object",
      "additionalProperties": false,
      "properties": {
        "regime_detected": {
          "type": "boolean"
        },
        "regime_type": {
          "type": "string",
          "enum": [
            "persistent_positive",
            "persistent_negative",
            "transitional",
            "low_conviction"
          ]
        },
        "positive_days": {
          "type": "integer"
        },
        "negative_days": {
          "type": "integer"
        },
        "avg_magnitude_billions": {
          "type": "number"
        },
        "sign_flips": {
          "type": "integer"
        },
        "persistence_pct": {
          "type": "number"
        },
        "confidence": {
          "type": "integer"
        },
        "reasoning": {
          "type": "string"
        }
      },
      "required": [
        "regime_detected",
        "regime_type",
        "positive_days",
        "negative_days",
        "avg_magnitude_billions",
        "sign_flips",
        "persistence_pct",
        "confidence",
        "reasoning"
      ]
    }
  }
}
\end{lstlisting}

\subsection{Request Body}
\label{app:prompt:request}

Each line of a batch input file is one request of the following form.

\begin{lstlisting}[style=appcode]
{
  "custom_id": "<arm>|<symbol>|<window_end>",
  "method": "POST",
  "url": "/v1/chat/completions",
  "body": {
    "model": "gpt-6-luna",
    "messages": [
      {
        "role": "user",
        "content": "<prompt text of Appendix A.2>"
      }
    ],
    "max_completion_tokens": 4000,
    "reasoning_effort": "medium",
    "response_format": {
      "type": "json_schema",
      "json_schema": "<schema of Appendix A.3>"
    }
  }
}
\end{lstlisting}

%=================================================================
\section{Series Definitions in SQL}
\label{app:series}

% src: 02_decisions_log.md D-08 (vendor cap 2026-10-09, cboe_delayed from 2026-10-12, irregular days 2026-07-29, 2026-08-03, 2026-09-14, series validity 2026-07-27..2026-10-09); README.md line 112. The verbatim audit statements are not stored as sql_series_*.sql in scripts/validation/electronics/sql (only gex_llm_regime_schema.sql); authors to place them there.
The daily series of Table~\ref{tab:series} are defined by the four read-only statements of Appendices~\ref{app:series:old} to~\ref{app:series:intraday}, which are reproduced verbatim from the data audit of 2026-10-28 on the project database; Appendix~\ref{app:series:derived} defines the derived and recovered series. The comments and the file names in the usage lines are those of the audit files. The audit calls the archive table ``frozen'' and its convention ``old''; Section~\ref{sec:methods:lineage} shows that this table is put-less from October 2023, so it is not a clean old convention (series key \ident{old_bs}). In the statements \texttt{:sym} is the symbol (SPY, QQQ or IWM) and \texttt{:y0}, \texttt{:y1} bound one calendar year, because the raw table is read one symbol-year per statement. Every statement stops at 2026-10-09, the last day before the vendor switch (Section~\ref{sec:methods:lineage}). The first two statements return the 30 daily values of every window. The pipeline stores the daily rows and cuts the windows in a second step (Section~\ref{sec:methods:windows}); unlike that step, the first statement drops a window when one of its 30 days of the \path{gex.options_daily_summary} calendar is missing from the archive table.

\subsection{Pre-Rebuild Archive Table (\texttt{old\_bs})}
\label{app:series:old}

\begin{lstlisting}[style=appcode]
-- SERIES (a) FROZEN old convention (read-only archive). usage: psql -v sym=SPY -f sql_series_a_frozen.sql
-- Calendar = NEW table calendar for the symbol; a 30-day window is valid only if all 30 consecutive NEW-calendar days exist in the frozen table.
-- (frozen max trading_date = 2026-10-02, so the last window ends 2026-10-02; IWM/QQQ have no rows 2020-12-03..2025-12-30)
WITH cal AS (SELECT trading_date, row_number() OVER (ORDER BY trading_date) AS rn
             FROM gex.options_daily_summary WHERE symbol = :'sym' AND trading_date <= DATE '2026-10-09'),
     f AS (SELECT c.rn, c.trading_date, z.total_gex, z.net_call_gex, z.net_put_gex, z.underlying_price
           FROM cal c JOIN backups.options_daily_summary_pre_rebuild_20261004 z
             ON z.symbol = :'sym' AND z.trading_date = c.trading_date)
SELECT e.trading_date AS window_end,
       array_agg(w.total_gex     ORDER BY w.rn) AS total_gex_30,
       array_agg(w.net_call_gex  ORDER BY w.rn) AS call_gex_30,
       array_agg(w.net_put_gex   ORDER BY w.rn) AS put_gex_30
FROM f e JOIN f w ON w.rn BETWEEN e.rn - 29 AND e.rn
GROUP BY e.trading_date, e.rn
HAVING count(*) = 30          -- hole-free window
ORDER BY e.rn;
\end{lstlisting}

\subsection{Rebuilt Daily Table (\texttt{new\_full})}
\label{app:series:full}

\begin{lstlisting}[style=appcode]
-- SERIES (b) NEW 0DTE-inclusive stored-gamma convention (gex.options_daily_summary, rebuilt 2026-10-08). usage: psql -v sym=SPY -f sql_series_b_new.sql
WITH s AS (SELECT trading_date, total_gex, net_call_gex, net_put_gex, underlying_price,
                  row_number() OVER (ORDER BY trading_date) AS rn
           FROM gex.options_daily_summary
           WHERE symbol = :'sym' AND trading_date <= DATE '2026-10-09')      -- alpha_vantage ends 2026-10-09; cboe_delayed starts 2026-10-12
SELECT e.trading_date AS window_end,
       array_agg(w.total_gex    ORDER BY w.rn) AS total_gex_30,
       array_agg(w.net_call_gex ORDER BY w.rn) AS call_gex_30,
       array_agg(w.net_put_gex  ORDER BY w.rn) AS put_gex_30
FROM s e JOIN s w ON w.rn BETWEEN e.rn - 29 AND e.rn
WHERE e.rn >= 30
GROUP BY e.trading_date, e.rn
HAVING count(*) = 30
ORDER BY e.rn;
\end{lstlisting}

\subsection{Raw Chains Without Same-Day Expiries (\texttt{new\_ex0dte})}
\label{app:series:ex0dte}

\begin{lstlisting}[style=appcode]
-- SERIES (c) NEW ex-0DTE from raw chains. ONE symbol-year per call (never wider). usage:
--   psql -v sym=SPY -v y0=2024-01-01 -v y1=2025-01-01 -f sql_series_c_ex0dte.sql      (loop y0/y1 over years; SPY/IWM 2008..2026, QQQ 2011..2026)
-- enable_seqscan=off is a SESSION planner setting (no writes): the default plan seq-scans the whole yearly partition (~4.5 s warm / 8.9 s cold per symbol-year);
-- the (symbol, trading_date) index/bitmap plan takes 0.5-1.8 s and returns identical numbers (max rel diff 5e-13).
SET enable_seqscan = off;
SELECT trading_date,
       count(*)                                                 AS n_contracts,
       count(*) FILTER (WHERE expiration = trading_date)        AS n_0dte,
       sum(CASE WHEN option_type = 'call' THEN g END)  FILTER (WHERE expiration > trading_date) AS call_gex_ex0dte,
      -sum(CASE WHEN option_type = 'put'  THEN g END)  FILTER (WHERE expiration > trading_date) AS put_gex_ex0dte,
       sum(CASE WHEN option_type = 'call' THEN g ELSE -g END) FILTER (WHERE expiration > trading_date) AS total_gex_ex0dte,
       sum(CASE WHEN option_type = 'call' THEN g ELSE -g END) FILTER (WHERE expiration = trading_date) AS total_gex_0dte,
       sum(CASE WHEN option_type = 'call' THEN g ELSE -g END)                                          AS total_gex_all  -- ties to gex.options_daily_summary.total_gex (<=2e-7 rel)
FROM (SELECT trading_date, option_type, expiration,
             gamma::float8 * open_interest::float8 * 100.0 * underlying_price::float8 * underlying_price::float8 * 0.01 AS g   -- float8: real*int would stay float4
      FROM gex.options_chains_partitioned
      WHERE symbol = :'sym' AND trading_date >= :'y0' AND trading_date < :'y1'
        AND trading_date <= DATE '2026-10-09'                     -- vendor cap: alpha_vantage ends 2026-10-09, cboe_delayed starts 2026-10-12
        AND gamma IS NOT NULL AND open_interest > 0 AND underlying_price > 0) t
GROUP BY trading_date
ORDER BY trading_date;
\end{lstlisting}

\subsection{Late-Session Intraday Snapshot (\texttt{intraday\_late})}
\label{app:series:intraday}

The statement is the one tested in the audit and returns Variant~1 of its comments, the \ident{market_close} slot, which equals the end-of-day row. The series of the study is Variant~2: the filter on the clock time is replaced by a filter on the snapshot type (\ident{snapshot_type} equal to \ident{final_rush}, 19:55:59 UTC), so that the snapshot is selected by type and not by clock time, and the three irregular days 2026-07-29, 2026-08-03 and 2026-09-14 are dropped when the series is built.

\begin{lstlisting}[style=appcode]
-- SERIES (d) INTRADAY late-session snapshot, SPY. Valid for 2026-07-27..2026-10-09 (EDT: 20:00 UTC = 16:00 ET; re-derive the cut after 2026-11-01 DST end).
-- Variant 1 'late' = last snapshot with timestamp < 20:01:00 UTC  -> the market_close slot (20:00:59), which EQUALS the daily summary row (not independent).
-- Variant 2 (independent of the EOD row): add  AND snapshot_type = 'final_rush'  (19:55:59 UTC).
SELECT DISTINCT ON (d)
       d AS trading_date, snapshot_timestamp, snapshot_type, total_gex, net_call_gex, net_put_gex, underlying_price, total_contracts
FROM (SELECT *, (snapshot_timestamp AT TIME ZONE 'America/New_York')::date AS d
      FROM gex.intraday_gex_summary
      WHERE symbol = 'SPY'
        AND snapshot_timestamp >= TIMESTAMPTZ '2026-07-27 00:00+00' AND snapshot_timestamp < TIMESTAMPTZ '2026-10-10 00:00+00'
        AND snapshot_type <> 'test'
        AND (snapshot_timestamp AT TIME ZONE 'UTC')::time < TIME '20:01:00') x
ORDER BY d, snapshot_timestamp DESC;
\end{lstlisting}

\subsection{Derived and Recovered Series}
\label{app:series:derived}

\textbf{Call leg only (\ident{call_only}).} The sum of the call terms \texttt{g} of the third statement over all expiries, same-day expiries included, with the put terms dropped. In the rebuilt summary table this is \ident{net_call_gex} (up to the float4 rounding of the stored column), which the second statement returns as \texttt{call\_gex\_30}.

\textbf{Scaled series (\ident{new_full_scaled}).} The \ident{new_full} values of QQQ and IWM multiplied by one fixed factor per symbol, which the following statement computes (tested read-only on 2026-10-29). The factor and the number of days are stored with the series.

\begin{lstlisting}[style=appcode]
-- new_full_scaled: factor of one symbol = median over 2020-01-01..2026-10-09 of the ratio of SPY's open-interest notional to the symbol's.
-- usage: psql -v sym=QQQ     (QQQ or IWM; the scaled values are the new_full values of the symbol times this factor)
WITH n AS (SELECT symbol, trading_date, total_weighted_oi::float8 * 100.0 * underlying_price::float8 AS notional
           FROM gex.options_daily_summary
           WHERE symbol IN ('SPY', :'sym') AND trading_date BETWEEN DATE '2020-01-01' AND DATE '2026-10-09')
SELECT percentile_cont(0.5) WITHIN GROUP (ORDER BY s.notional / q.notional) AS factor,
       count(*) AS n_days
FROM n s JOIN n q ON q.trading_date = s.trading_date AND s.symbol = 'SPY' AND q.symbol = :'sym'
WHERE q.notional > 0;
\end{lstlisting}

\textbf{Recovered input series.} The series \ident{jan2026_phase4a}, \ident{p3_2024} and \ident{p4_2020} (and the dated twins of the last two) are not defined by SQL. They are the daily values parsed from the data lines of the surviving batch input files \path{phase4a_postgresql_20260104_184832.jsonl}, \path{phase3_full_2024_baseline_v2.jsonl} and \path{batch_regime_20251120_104630.jsonl} (folder \path{reports/validation/paper2_regime_windows/batch_jobs/}); overlapping windows are checked for consistent values (Section~\ref{sec:results:inputs}).

%=================================================================
\section{Per-Year Tables}
\label{app:tables}

Tables~\ref{tab:app:series}--\ref{tab:ledger} give the numbers behind the figures of Section~\ref{sec:results} for every year, taken from the files that \path{scripts/validation/electronics/05_report.py} writes to \path{docs/papers/electronics/results/}. The years are the calendar years of the window end, and rates are percentages of parsed windows. Wilson and quarter-block bootstrap intervals accompany every rate in \path{detection_by_year.csv} and \path{stats.yaml}; the Wilson intervals treat overlapping windows as independent and are optimistic.

\begin{table}[H]
\centering
\caption{Mean absolute daily GEX (USD billion) by series and year, SPY. Each mean is taken over the days that the series has in common with \ident{new_full}; the \ident{new_full} column and the day count use the days shared with \ident{new_ex0dte}. A dash marks a year outside the series.}
% src: series_comparison.csv scope=year symbol=SPY; mean_abs_a_bn / mean_abs_b_bn and n_days of the pairs jan2026_phase4a_vs_new_full, call_only_vs_jan2026_phase4a, old_bs_vs_new_full, new_full_vs_new_ex0dte
\label{tab:app:series}
\footnotesize
\begin{tabular}{@{}lrrrrrr@{}}
\toprule
Year & Days & \ident{jan2026_phase4a} & \ident{call_only} & \ident{old_bs} & \ident{new_ex0dte} & \ident{new_full} \\
\midrule
2008 & 253 & -- & -- & 1.39 & 1.59 & 1.60 \\
2009 & 250 & -- & -- & 0.98 & 0.88 & 0.88 \\
2010 & 251 & -- & -- & 1.01 & 1.14 & 1.16 \\
2011 & 250 & -- & -- & 1.79 & 1.99 & 2.02 \\
2012 & 250 & -- & -- & 1.54 & 1.71 & 1.78 \\
2013 & 252 & -- & -- & 2.34 & 2.29 & 2.51 \\
2014 & 252 & -- & -- & 3.03 & 3.07 & 3.36 \\
2015 & 252 & -- & -- & 2.69 & 2.91 & 3.37 \\
2016 & 250 & -- & -- & 3.99 & 3.76 & 4.23 \\
2017 & 240 & -- & -- & 6.43 & 5.82 & 6.53 \\
2018 & 249 & -- & -- & 4.79 & 4.63 & 5.32 \\
2019 & 251 & -- & -- & 3.62 & 3.39 & 4.11 \\
2020 & 252 & 3.24 & 10.51 & 2.76 & 2.57 & 3.24 \\
2021 & 252 & 4.97 & 17.13 & 3.98 & 3.50 & 4.99 \\
2022 & 251 & 5.48 & 9.27 & 4.42 & 4.62 & 5.46 \\
2023 & 250 & 10.62 & 17.04 & 8.39 & 4.59 & 6.10 \\
2024 & 252 & 20.19 & 20.19 & 16.19 & 4.13 & 6.46 \\
2025 & 250 & 19.21 & 19.21 & 14.71 & 4.55 & 7.13 \\
2026 & 193 & -- & -- & 6.00 & 5.76 & 8.37 \\
\bottomrule
\end{tabular}
\end{table}

\begin{table}[H]
\centering
\caption{Detection rate of the deterministic classifier by series and window-end year, SPY, in percent of windows. Windows are the windows of \ident{new_full}. A dash marks a year without windows in the series.}
% src: detection_by_year.csv label_source=classifier scope=year symbol=SPY, column n (windows of new_full) and rate_pct per series_key
\label{tab:app:replay}
\footnotesize
\begin{tabular}{@{}lrrrrrr@{}}
\toprule
Year & Windows & \ident{jan2026_phase4a} & \ident{call_only} & \ident{old_bs} & \ident{new_ex0dte} & \ident{new_full} \\
\midrule
2008 & 224 & -- & -- & 0.0 & 0.0 & 0.0 \\
2009 & 250 & -- & -- & 0.0 & 0.0 & 0.0 \\
2010 & 251 & -- & -- & 0.0 & 0.0 & 0.0 \\
2011 & 250 & -- & -- & 0.0 & 0.0 & 0.0 \\
2012 & 250 & -- & -- & 0.0 & 0.0 & 0.0 \\
2013 & 252 & -- & -- & 0.0 & 0.0 & 0.0 \\
2014 & 252 & -- & -- & 13.5 & 8.3 & 9.1 \\
2015 & 252 & -- & -- & 0.4 & 1.6 & 7.9 \\
2016 & 250 & -- & -- & 26.4 & 22.0 & 28.0 \\
2017 & 240 & -- & -- & 38.8 & 19.6 & 32.1 \\
2018 & 249 & -- & -- & 32.1 & 23.7 & 34.5 \\
2019 & 251 & -- & -- & 7.2 & 3.2 & 10.0 \\
2020 & 252 & 12.1 & 90.6 & 16.7 & 10.7 & 22.2 \\
2021 & 252 & 4.0 & 100.0 & 3.6 & 0.0 & 3.6 \\
2022 & 251 & 32.3 & 100.0 & 26.3 & 36.7 & 40.6 \\
2023 & 250 & 19.3 & 100.0 & 22.5 & 25.2 & 22.4 \\
2024 & 252 & 100.0 & 100.0 & 100.0 & 13.1 & 10.3 \\
2025 & 250 & 100.0 & 100.0 & 100.0 & 16.4 & 13.6 \\
2026 & 193 & -- & 100.0 & 23.5 & 21.2 & 11.4 \\
\bottomrule
\end{tabular}
\end{table}

\begin{table}[H]
\centering
\caption{Detection rate of gpt-6-luna by arm and window-end year, in percent of parsed windows (a dash marks a year without windows in the arm). The columns are the arms \ident{full_spy}, \ident{ex0dte_spy}, \ident{full_qqq}, \ident{full_qqq_scaled}, \ident{ex0dte_qqq}, \ident{full_iwm}, \ident{full_iwm_scaled} and \ident{ex0dte_iwm} of Table~\ref{tab:arms} (full = \ident{new_full}, scaled = \ident{new_full_scaled}, ex-0DTE = \ident{new_ex0dte}); the bridge, intraday and Phase 3/4 arms are reported in Sections~\ref{sec:results:bridge}, \ref{sec:results:intraday} and~\ref{sec:results:dated}.}
\label{tab:app:llm}
% src: detection_by_year.csv label_source=llm arm_role=primary scope=year, column rate_pct per arm
\footnotesize
\begin{tabular}{@{}lrrrrrrrr@{}}
\toprule
 & \multicolumn{2}{c}{SPY} & \multicolumn{3}{c}{QQQ} & \multicolumn{3}{c}{IWM} \\
\cmidrule(lr){2-3}\cmidrule(lr){4-6}\cmidrule(lr){7-9}
Year & full & ex-0DTE & full & scaled & ex-0DTE & full & scaled & ex-0DTE \\
\midrule
2008 & 0.0 & 0.0 & -- & -- & -- & 0.0 & -- & -- \\
2009 & 0.0 & 0.0 & -- & -- & -- & 0.0 & -- & -- \\
2010 & 0.0 & 0.0 & -- & -- & -- & 0.0 & -- & -- \\
2011 & 0.0 & 0.0 & 0.0 & -- & -- & 0.0 & -- & -- \\
2012 & 0.0 & 0.0 & 0.0 & -- & -- & 0.0 & -- & -- \\
2013 & 0.0 & 0.0 & 0.0 & -- & -- & 0.0 & -- & -- \\
2014 & 9.1 & 8.3 & 0.0 & -- & -- & 0.0 & -- & -- \\
2015 & 7.9 & 1.6 & 0.0 & -- & -- & 0.0 & -- & -- \\
2016 & 27.6 & 22.0 & 0.0 & -- & -- & 0.0 & -- & -- \\
2017 & 31.7 & 19.2 & 0.0 & -- & -- & 0.0 & -- & -- \\
2018 & 34.9 & 22.9 & 0.0 & -- & -- & 0.0 & -- & -- \\
2019 & 10.0 & 3.2 & 0.0 & -- & -- & 0.0 & -- & -- \\
2020 & 22.2 & 10.7 & 0.0 & 0.0 & 0.0 & 0.0 & 9.9 & 0.0 \\
2021 & 3.6 & 0.0 & 0.0 & 23.8 & 0.0 & 0.0 & 65.9 & 0.0 \\
2022 & 41.8 & 37.1 & 0.0 & 61.0 & 0.0 & 0.0 & 100.0 & 0.0 \\
2023 & 23.2 & 26.0 & 0.0 & 17.6 & 0.0 & 0.0 & 79.2 & 0.0 \\
2024 & 11.1 & 14.3 & 0.0 & 8.4 & 0.0 & 0.0 & 20.6 & 0.0 \\
2025 & 14.0 & 16.4 & 0.0 & 20.4 & 0.0 & 0.0 & 51.2 & 0.0 \\
2026 & 11.4 & 21.2 & 0.0 & 16.1 & 0.0 & 0.0 & 45.1 & 0.0 \\
\bottomrule
\end{tabular}
\end{table}

% src: cost_ledger.csv, rows with counts_toward_budget = True (column recomputed_cost_usd; it equals est_cost_usd in every row); total equals cost_ledger.yaml budget.spent_usd (9.700576); backfill rows budget_class = backfill
\begin{table}[H]
\centering
\caption{Run ledger: one row per Batch job, with the token usage and the cost priced at Batch rates. The identifier has the form \ident{<arm>_<model>_<mode>_<UTC timestamp>_c<chunk>}. Only runs paid in this study enter the total; the synchronous pilot (\ident{sync}) is included in it. The three backfilled o4-mini runs are listed below it, as stored answers of earlier runs that are not spend of this study.}
\label{tab:ledger}
\scriptsize
\setlength{\tabcolsep}{3pt}
\begin{tabular}{@{}p{4.9cm}rrrrl@{}}
\toprule
Run identifier & Requests & Prompt tokens & Completion tokens & Cost (USD) & Status \\
\midrule

\path{bridge_jan_gpt-6-luna_batch_20261029T203742Z_c0} & 1{,}476 & 2{,}823{,}588 & 1{,}086{,}712 & 0.41 & ingested \\
\path{bridge_jan_gpt-6-luna_batch_20261029T203742Z_c1} & 1{,}476 & 2{,}823{,}588 & 1{,}082{,}280 & 0.41 & ingested \\
\path{full_spy_gpt-6-luna_batch_20261029T203745Z_c0} & 1{,}227 & 2{,}347{,}251 & 871{,}277 & 0.34 & ingested \\
\path{full_spy_gpt-6-luna_batch_20261029T203745Z_c1} & 1{,}226 & 2{,}345{,}338 & 969{,}842 & 0.36 & ingested \\
\path{full_spy_gpt-6-luna_batch_20261029T203745Z_c2} & 1{,}226 & 2{,}345{,}338 & 1{,}004{,}850 & 0.37 & ingested \\
\path{full_spy_gpt-6-luna_batch_20261029T203745Z_c3} & 1{,}226 & 2{,}345{,}338 & 1{,}024{,}540 & 0.37 & ingested \\
\path{ex0dte_spy_gpt-6-luna_batch_20261029T203749Z_c0} & 1{,}227 & 2{,}347{,}251 & 839{,}753 & 0.33 & ingested \\
\path{ex0dte_spy_gpt-6-luna_batch_20261029T203749Z_c1} & 1{,}226 & 2{,}345{,}338 & 968{,}965 & 0.36 & ingested \\
\path{ex0dte_spy_gpt-6-luna_batch_20261029T203749Z_c2} & 1{,}226 & 2{,}345{,}338 & 951{,}505 & 0.36 & ingested \\
\path{ex0dte_spy_gpt-6-luna_batch_20261029T203749Z_c3} & 1{,}226 & 2{,}345{,}338 & 982{,}560 & 0.36 & ingested \\
\path{p34_obf_2024_gpt-6-luna_batch_20261029T203752Z_c0} & 234 & 447{,}642 & 156{,}636 & 0.06 & ingested \\
\path{p34_dated_2024_gpt-6-luna_batch_20261029T203754Z_c0} & 234 & 461{,}682 & 154{,}330 & 0.06 & ingested \\
\path{p34_obf_2020_gpt-6-luna_batch_20261029T203757Z_c0} & 234 & 447{,}642 & 132{,}619 & 0.06 & ingested \\
\path{p34_dated_2020_gpt-6-luna_batch_20261029T203759Z_c0} & 234 & 461{,}682 & 130{,}805 & 0.06 & ingested \\
\path{full_qqq_gpt-6-luna_batch_20261029T204029Z_c0} & 1{,}354 & 2{,}590{,}202 & 1{,}012{,}074 & 0.38 & ingested \\
\path{full_qqq_gpt-6-luna_batch_20261029T204029Z_c1} & 1{,}353 & 2{,}588{,}289 & 1{,}004{,}705 & 0.38 & ingested \\
\path{full_qqq_gpt-6-luna_batch_20261029T204029Z_c2} & 1{,}353 & 2{,}588{,}289 & 1{,}057{,}334 & 0.39 & ingested \\
\path{full_iwm_gpt-6-luna_batch_20261029T204033Z_c0} & 1{,}228 & 2{,}349{,}164 & 762{,}229 & 0.31 & ingested \\
\bottomrule
\end{tabular}
\end{table}

\begin{table}[H]
\centering
\caption{Run ledger (continued): remaining Batch jobs, the total of this study and the backfilled runs.}
\label{tab:ledger2}
\scriptsize
\setlength{\tabcolsep}{3pt}
\begin{tabular}{@{}p{4.9cm}rrrrl@{}}
\toprule
Run identifier & Requests & Prompt tokens & Completion tokens & Cost (USD) & Status \\
\midrule
\path{full_iwm_gpt-6-luna_batch_20261029T204033Z_c1} & 1{,}228 & 2{,}349{,}164 & 719{,}759 & 0.30 & ingested \\
\path{full_iwm_gpt-6-luna_batch_20261029T204033Z_c2} & 1{,}228 & 2{,}349{,}164 & 826{,}909 & 0.32 & ingested \\
\path{full_iwm_gpt-6-luna_batch_20261029T204033Z_c3} & 1{,}228 & 2{,}349{,}164 & 855{,}981 & 0.33 & ingested \\
\path{full_qqq_scaled_gpt-6-luna_batch_20261029T204036Z_c1} & 877 & 1{,}677{,}701 & 750{,}362 & 0.27 & ingested \\
\path{full_qqq_scaled_gpt-6-luna_batch_20261029T204036Z_c0} & 877 & 1{,}677{,}701 & 694{,}967 & 0.26 & ingested \\
\path{full_iwm_scaled_gpt-6-luna_batch_20261029T204039Z_c0} & 878 & 1{,}679{,}614 & 614{,}557 & 0.24 & ingested \\
\path{full_iwm_scaled_gpt-6-luna_batch_20261029T204039Z_c1} & 877 & 1{,}677{,}701 & 695{,}367 & 0.26 & ingested \\
\path{ex0dte_qqq_gpt-6-luna_batch_20261030T010447Z_c1} & 892 & 1{,}706{,}396 & 680{,}360 & 0.26 & ingested \\
\path{ex0dte_qqq_gpt-6-luna_batch_20261030T010447Z_c0} & 892 & 1{,}706{,}396 & 622{,}593 & 0.24 & ingested \\
\path{ex0dte_iwm_gpt-6-luna_batch_20261030T010451Z_c0} & 893 & 1{,}708{,}309 & 560{,}829 & 0.23 & ingested \\
\path{ex0dte_iwm_gpt-6-luna_batch_20261030T010451Z_c1} & 892 & 1{,}706{,}396 & 616{,}040 & 0.24 & ingested \\
\path{intraday_spy_gpt-6-luna_batch_20261029T204042Z_c0} & 54 & 103{,}302 & 45{,}525 & 0.02 & ingested \\
\path{bridge_old_gpt-6-luna_batch_20261030T010455Z_c0} & 1{,}218 & 2{,}330{,}034 & 813{,}090 & 0.32 & ingested \\
\path{bridge_old_gpt-6-luna_batch_20261030T010455Z_c1} & 1{,}218 & 2{,}330{,}034 & 970{,}974 & 0.36 & ingested \\
\path{bridge_old_gpt-6-luna_batch_20261030T010455Z_c2} & 1{,}217 & 2{,}328{,}121 & 953{,}927 & 0.35 & ingested \\
\path{bridge_old_gpt-6-luna_batch_20261030T010455Z_c3} & 1{,}217 & 2{,}328{,}121 & 830{,}288 & 0.32 & ingested \\
\path{pilot_bridge_jan_gpt-6-luna_sync_20261029T203454Z_c0} & 40 & 76{,}520 & 28{,}015 & 0.02 & ingested \\
\midrule
Total (this study) & 34{,}712 & 66{,}432{,}136 & 25{,}472{,}559 & 9.70 & \\
\midrule
\multicolumn{6}{@{}l}{Backfilled o4-mini runs (not spend of this study)} \\
\path{r1_o4mini_phase4a_20260104} & 1{,}476 & 2{,}655{,}324 & 2{,}928{,}992 & 7.90 & backfill \\
\path{r1_p3_o4mini_20251120} & 223 & 414{,}557 & 405{,}127 & 1.12 & backfill \\
\path{r1_p4_o4mini_20251120} & 223 & 414{,}557 & 325{,}436 & 0.94 & backfill \\
\bottomrule
\end{tabular}
\end{table}

%=================================================================
\section{Data Availability Details}
\label{app:data}

% src: data/backups/gex_llm/condensed/ (row counts read from the files, 2026-10-31); series_index.csv; scripts/validation/electronics/README.md (storage section); 02_decisions_log.md D-14, D-17
The repository holds the condensed material and the results. The condensed files are in \path{data/backups/gex_llm/condensed/}.
\begin{itemize}[leftmargin=*,labelsep=5.5mm]
\item \path{detections.csv.gz}: 36{,}634 rows, one per stored model answer, with the verdict, the reported metrics and the parse status but without the reasoning or the raw text. The key is \ident{arm}, \ident{series_key}, \ident{symbol}, \ident{window_end}, \ident{model}, \ident{prompt_version}.
\item \path{windows.csv.gz}: 38{,}501 rows, one per 30-day window, with the labels, the sha256 of the input as sent and the verdict of the deterministic rule but without the sequences.
\item \path{series_index.csv}: 16 rows, one per series and symbol, with the row count, the first and last date, the source table and the hash of the builder.
\end{itemize}
Every number of the paper can be reproduced from the first two files (join on \ident{series_key}, \ident{symbol} and \ident{window_end}) and the files of \path{docs/papers/electronics/results/}. The provenance of every Batch job (batch and file identifiers, sha256 of the input file, token counts and cost) is in \path{data/backups/gex_llm/tables/runs/} and in the \path{manifest.json} and \path{batch_<k>.json} files of each run folder. The raw request and answer files (\path{input_<k>.jsonl}, \path{output_<k>.jsonl}) and the full tables are kept on the authors' machine and in the database; they are not part of the repository.

The full tables are loaded into four tables and one view of the schema \ident{backups} of the research database: \ident{gex_llm_regime_series} (the number shown to the model, one row per series, symbol and trading date), \ident{gex_llm_regime_windows} (the 30-day input as sent, its sha256 and the verdict of the rule), \ident{gex_llm_regime_runs} (one row per Batch job, pilot or backfilled run) and \ident{gex_llm_regime_detections} (one row per model answer), with the view \ident{gex_llm_regime_rates_summary} (yearly rate of the model next to the rate of the rule, by arm). The parquet copies of these four tables are the release format named in the Data Availability Statement. The options data are subject to the vendor's licence and are not redistributed.
